\title{[0pt][l]{-1ex} SeaLLMs - Large Language Models for Southeast Asia}

\begin{abstract}
Despite the remarkable achievements of large language models (LLMs) in various tasks, there remains a linguistic bias that favors high-resource languages, such as English, often at the expense of low-resource and regional languages. To address this imbalance, we introduce SeaLLMs, an innovative series of language models that specifically focuses on Southeast Asian (SEA) languages. SeaLLMs are built upon popular English-centric models through continued pre-training with an extended vocabulary, specialized instruction and alignment tuning to better capture the intricacies of regional languages. This allows them to respect and reflect local cultural norms, customs, stylistic preferences, and legal considerations. Our comprehensive evaluation demonstrates that SeaLLM models exhibit superior performance across a wide spectrum of linguistic tasks and assistant-style instruction-following capabilities relative to comparable open-source models. Moreover, they outperform ChatGPT-3.5 in non-Latin languages, such as Thai, Khmer, Lao, and Burmese, by large margins while remaining lightweight and cost-effective to operate.
\end{abstract}

\section{Introduction}\label{sec:intro}

The advent of large language models (LLMs) has radically transformed the field of natural language processing, demonstrating remarkable abilities in text generation, comprehension, and decision-making tasks \cite{gpt3_brown2020language,chatgpt,gpt4,llama_touvron2023,llama2touvron2023llama,lambda_google_thoppilan2022lamda,mistral7b_jiang2023mistral,polylm_wei2023polylm,qwen_bai2023}. While the proficiencies of these models are extraordinary, the majority of existing LLMs embody a linguistic hierarchy overwhelmingly dominated by English \cite{mega, multilingual-eval, m3exam}. This dominance undermines the multilingual capability of such models, with particularly prejudicial outcomes for lower-resource and regional languages, where data scarcity and tokenization challenges lead to disproportionately poor model performance. This linguistic disparity not only impedes access to state-of-the-art AI technologies for non-English-speaking populations but also risks cultural homogenization and the loss of linguistic diversity. While hyper-polyglot models exist \citep{bloom_scao2022bloom,bloomz_muennighoff2022crosslingual,polylm_wei2023polylm}, they may pay a high cost for high-resource language performance while lacking in multilingual instruction-following abilities.

Recognizing the urgent need to democratize AI and empower linguistically diverse regions, we introduce SeaLLMs\footnote{\href{https://github.com/DAMO-NLP-SG/SeaLLMs}{https://github.com/DAMO-NLP-SG/SeaLLMs}}, a suite of specialized language models optimized for Southeast Asian languages\footnote{English (Eng), Chinese (Zho), Indonesian (Ind), Vietnamese (Vie), Thai (Tha), Khmer (Khm), Lao, Malay (Msa), Burmese (Mya) and Tagalog (Tgl)}. These languages, while rich and diverse, often lack the extensive dataset support available for more widely spoken languages, resulting in a stark performance gap in existing LLM applications. 

As a long-term continuous effort, as of this writing, SeaLLMs come in three versions (v1, v2, v2.5). SeaLLM-13B-v1, which was pre-trained from Llama-2-13B, eclipses the performance of most available open-source LLMs in a comprehensive array of tasks including world knowledge assessments, language comprehension, and generative capabilities in SEA languages. For English and alike, SeaLLMs do not only preserve, but also demonstrate enhanced performance in tasks that were part of the original Llama training set. When evaluated on multilingual instruction-following tasks with GPT-4 as a judge \citep{mt_bench_zheng2023judging}, SeaLLM-13B-v1 outperforms ChatGPT-3.5 by large margins in less-represented languages such as Khmer, Lao or Burmese. Meanwhile, SeaLLM-7B-v2, which was pre-trained from Mistral-7B \citep{mistral7b_jiang2023mistral}, demonstrates better performances in math and commonsense reasoning than comparable baselines, surpassing ChatGPT-3.5 in reasoning for common SEA languages, while being much smaller in sizes. Later, SeaLLM-7B-v2.5, which was further pre-trained from Gemma-7B \citep{gemma_team2024}, shows significant improvements in SEA languages over SeaLLM-7B-v2.

\ref{fig:seallm_train_process} illustrates the four-stage training process of SeaLLMs. In the first stage, detailed in \ref{sec:pretrain:process}, we conduct continuous pre-training from the foundational models \citep{llama2touvron2023llama,mistral7b_jiang2023mistral} with an extended vocabulary tailored for SEA languages. Next, we fine-tune the model in a novel hybrid paradigm with a mixture of multilingual pre-training data and English-dominant instruction fine-tuning data (\ref{sec:sft:pretrain_sft}). The following stage subsequently fine-tunes the model on a balanced and custom-built multilingual SFT dataset. Finally, we conduct self-preferencing alignment optimization using the SeaLLM model itself, without relying on human annotators or more powerful LLMs \citep{gpt4}.

\section{Related Work}\label{sec:rel_works}

\makeatletter

\pgfplotsset{
  col chart axis style/.style={
    width=\columnwidth,
    height=0.15\textheight,
    xtick=data,
    enlarge y limits=false, enlarge x limits=0.05, }, col bar axis style/.style={ col chart axis style, enlarge y limits=false, enlarge x limits=0.05, ymin=40, ymax=65, ybar=0.1em, bar width=9pt, legend style={
      font=\footnotesize,

      legend columns=2,

      at={(0.15, 0.95)},  

      anchor=north, /tikz/every even column/.append style={column sep=0.5cm}, }, }, }

Large language models (LLMs) display outstanding capabilities because they are pre-trained on massive amounts of internet text data \citep{gpt3_brown2020language,palm_chowdhery2022palm,bloom_scao2022bloom,llama_touvron2023,llama2_touvron2023llama}. Without any gradient update, base LLMs are able to perform in-context learning by simply observing a list of high-quality input-output exemplars \citep{gpt3_brown2020language,llm_do_icl_differently_wei2023larger}. This technique works across many tasks that involve language understanding, reasoning and generation \citep{gpt3_brown2020language,cot_wei2022chain,llm_are_mling_cot_shi2022language}. Much research has been conducted to understand in-context learning. Some suggest that the models secretly perform gradient descent on the exemplars \citep{icl_gradient_descen_dai2022can}, while others demonstrate that most of the knowledge is learned during pre-training, and the exemplars are only to provide evidence for the model to \textit{locate} the intended task via a Bayesian inference process \citep{icl_bayesian_infer_xie2021explanation,rethink_demonstration_min2022rethinking,lima_zhou2023lima}.

Most LLMs are trained with multilingual corpora \citep{ccnet_wenzek-etal-2020-ccnet}, even if these make up a tiny fraction of the large English corpora \citep{gpt2_radford2019language_gpt2,gpt3_brown2020language}. Despite that, LLMs still exhibit strong multilingual capabilities with high-resource languages like French, German and Chinese, often with the help of English-pivoting using supervised translation systems \citep{llm_are_mling_cot_shi2022language} or prompting the model to firstly generate intermediate English context \citep{xlt_not_all_created_equal_huang2023not}. 

BLOOM \citep{bloom_scao2022bloom} is one of the LLMs trained with the most number of languages, whose ROOTS corpus consists of data from 46 languages \citep{roots_corpus}. The ROOTS corpus includes 34 Indic and African languages regarded as low-resource, with each language having a pre-training coverage of less than 1\% in Hindi for the Indic group, to $2e^{-5}$\% in Tumbuka for the African group, as shown in \ref{fig:bloom_lang_coverage} in the Appendix. 

Exclusively only for unsupervised translation tasks, our linguistically-diverse prompting (LDP) strategy is also an English-pivoting method, but it is different from other cross-lingual counterparts \citep{llm_are_mling_cot_shi2022language,xlt_not_all_created_equal_huang2023not} in that while others only pivot inputs to English intermediates, we use in-context pairs between English and a diverse set of high-resource languages to promote the intended task in the target low-resource language. For other intra-lingual tasks like instruction following, where input and output are both expected to be in the same language, our LDP helps prevent English-tuned models \citet{llama2_touvron2023llama} from responding with English answer given a non-English query.

Part of our work also intersects with unsupervised multilingual machine translation (UMT), where back-translation is proven to be effective \citep{understanding_backtranslation_scale,lample2018phrase_unsup,conneau2019cross_xlm,mbart_liu2020multilingual,refine_lowres_umt_nguyen2022}, along with other methods \citep{criss2020,swavumt_nguyen2022contrastive}. English-pivoting is also prominent in the realm of machine translation, where training models on high-resource En$\leftrightarrow$X bitext improves lower-resource En$\leftrightarrow$Y tasks \citep{multilingual_view_garcia-etal-2020,harness-multilingual-garcia-etal-2021}. 

Analyses of machine translation using LLMs have also been done. \citet{howgood_gpt_mt_hendy2023good} show that GPT models can perform competitively alongside the best MT models. \citet{mmt_llm_analysis_zhu2023multilingual} focus on optimizing supervised exemplars selection strategies, while \citet{icl_maintain_coherence_sia2023context} discover that using specific coherent prompts for each input helps improve performance. Nonetheless, such studies only consider supervised instruction-tuned models \citep{instructgpt_ouyang2022training,bloomz_muennighoff2022crosslingual}, which may risk test-set contamination \citep{bloomz_muennighoff2022crosslingual}. Thus, there is still limited research involving low-resource languages in completely zero-shot setups. As such, since low-resource languages may not enjoy the privilege of having large unlabeled data to conduct searching, only random selection is used in this study, while optimal exemplar selection is out of scope.

\section{Pre-training}\label{sec:pretraining}

\subsection{Pre-training Data}\label{sec:pretrain:data}

The pre-training data comprises a heterogeneous assortment of documents sourced from several publicly accessible repositories~\citep{suarez2019asynchronous,mc4_2019t5,together2023redpajama,wikidump}. Specifically, during the creation of the pre-training data, we include web-based corpora such as Common Crawl \citep{ccnet_wenzek-etal-2020-ccnet}, journalistic content such as CC-News, text corpora with expertly-curated knowledge such as Wikipedia \citep{wikidump}, and some scholarly publications. After collecting the data, we employ a language identifier \citep{fasttext} to retain the documents for the major languages in Southeast Asia, namely Thai, Vietnamese, Indonesian, Chinese, Khmer, Lao, Malay, Burmese, and Tagalog, and discard the remaining ones. Subsequent stages of data refinement involve the multiple modules dedicated to data cleansing and content filtration. We blend such data with the highest quality English data from RedPajama subset \citep{together2023redpajama} in more balanced ratios, as we found that such English data are useful to preserve the original learnt knowledge.

\subsection{Vocabulary Expansion}\label{sec:pretrain:vocab_extend}

\ref{tab:language_ratios} describes how expensive it is to process an under-represented non-Latin language. For example, encoding a single sentence in Thai requires 4.3 times more tokens than its English equivalent. The reason for this is that most English language models employ a BPE tokenizer \citep{bpe_sennrich2015neural} that inefficiently segments texts from non-Latin scripts into disproportionately lengthy byte sequences, which inadequately represent the underlying semantic content, resulting in diminished model performance \citep{democratize_llm_ldp_nguyen2023}. 

To that end, we propose a novel vocabulary expansion technique, as formally described in \ref{algo:vocab_extend} in the Appendix. This technique involves recursively merging whole-word and sub-word token pieces of a new language from a highly multilingual target tokenizer (\emph{i.e.,}\ the NLLB tokenizer \citep{nllb_costa2022no_flores200}), to the existing LLM tokenizer. This new set of retrieved tokens are then pruned to remove rarely appearing and low-quality tokens before being added to the final SeaLLM tokenizer. 

\ref{tab:language_ratios} demonstrates the efficiency of the new vocabulary. The compression ratio for Thai text has markedly improved from 4.29 to 1.57, signifying a 2.7-fold increase in the length of Thai text that can be encoded within the same context constraints. At the same time, the compression of English text has experienced a negligible reduction of 0.3\%, thus maintaining its tokenization effectiveness.

We applied our vocabulary expansion for SeaLLM v1 and v2 with Llama-2 and Mistral-7B as backbones due to their limit 32K-token vocabulary. However, we did not extend the tokenizer for SeaLLM-7B-v2.5, which inherits a sufficiently large 250K-token vocabulary from Gemma-7B.

\subsection{Pre-training Process}\label{sec:pretrain:process}

We organize our pre-training dataset based on the language of the content and the quality of the data, as mentioned in \ref{sec:pretrain:data}. We setup a separate stream of data for each language, and dynamically control and balance the sampling ratio of each language. We pack multilingual documents into a single sequence up to the maximum context length. During the last steps of pre-training, we re-feed the model with more high quality data, which it has previously seen, to readjust the model's learning focus back towards the high-quality data, improving the model's performance.

\section{Supervised Fine-tuning (SFT)}\label{sec:sft}

\subsection{Supervised Fine-tuning Data}\label{sec:sft:data}

Our supervised finetuning (SFT) data consists of many categories, including text understanding and processing, math and logical reasoning, user-centric instruction-following, and natural dialog data. As most public and open-source SFT data are English-only~\citep{flan_collection_longpre2023,openorca,orca_mukherjee2023orca,platypus2023}, various techniques were implemented to enhance the multilingual aspect of the model. These include sourcing natural data from local websites in natural settings, selectively translating from English data, employing self-instruction, and using advanced prompting techniques~\citep{self_instruct_wang2022self,self_refine_madaan2023self,democratize_llm_ldp_nguyen2023}. As those synthetically generated data may remain incorrect or low-quality, native speakers\footnote{Hired by our organization, they are not co-authors.} were then engaged to further verify, filter, and edit such synthetic responses to finalize the SFT dataset. We find that engaging the annotators to verify and modify model-generated responses is more efficient than having them write responses from scratch. Safety-related data also played a crucial role in fine-tuning SeaLLMs. We manually collected and prepared country-relevant safety data, which covered a broad range of culturally and legally sensitive topics in each of these countries. This was necessary as such topics are often overlooked or may even conflict with open-source English-centric safety data \cite{multilingual_jailbreak_deng2023}. 

For SeaLLM-7B-v2 and SeaLLM-7B-v2.5, we incorporate significantly more SFT data relating to math and commonsense reasoning. Such data is synthetically generated with SeaLLM-13B-v1, as well as strong English models \citep{jiang2024mixtral,qwen_bai2023} using a combination of few-shot paraphrasing and translation techniques \citep{yu2023metamath}.

\subsection{Supervised Fine-tuning}\label{sec:sft:pretrain_sft}

\paragraph{Pre-train and SFT Hybrid.} As our SFT data is still significantly English due to contributions of open-source data, directly conducting SFT on it may overshadow the smaller SEA language datasets. Therefore, we propose incorporating an additional step prior to complete fine-tuning, namely Pre-train \& SFT Hybrid. In this step, the model is further trained on a combination of the pre-training corpus and a large portion of English SFT data, leaving the remaining and more balanced amount of English SFT data to the next stage. During this hybrid stage, the model processes both general pre-training content and instruction-following examples. We mask the source side of the instruction or supervised data to prevent the model from overfitting to the training examples and to reduce the risk of it simply memorizing the input data instead of learning the more generalized ability to follow instructions.

\paragraph{Supervised Fine-tuning.} We conduct supervised fine-tuning by compiling instructions from a variety of sources explained in \ref{sec:sft:data}, combining them at random into a single, consolidated sequence to maximize efficiency. To enhance the multi-turn conversation capability, in the later stage of fine-tuning, we further artificially create multi-turn conversations by randomly joining several single-turn instructions together. 

\subsection{Self-Preferencing Optimization}\label{sec:dpo}

Alignment from human feedback preference has been key to the success of many AI-assistant language models \citep{rlhf_stiennon2020learning,llama2touvron2023llama,dpo_rafailov2023direct,instructgpt_ouyang2022training}. To save the cost of human preference annotation work, some have sought to use powerful LLMs like GPT-4 \citep{gpt4} to play the part of a preference data generator \citep{zephyr7b_tunstall2023}. However, that may not even be feasible for low-resource non-Latin languages because of the unfavorable tokenization of ChatGPT as explained in \ref{sec:pretrain:vocab_extend}. In other words, even short prompts would exceed their context-length and the API-call costs would explode by up to 17 times.

Therefore, we use our own SeaLLM SFT models to generate preference data by asking it to indicate its preference between two of its own responses, given a question based on certain human-written criteria. To eliminate position bias, we swap the order of the responses and remove samples with inconsistent preference. The data is later used to employ direct preference optimization \citep{dpo_rafailov2023direct} to significantly improve the model abilities as an assistant. As such, unlike other works \citep{orca_mukherjee2023orca,zephyr7b_tunstall2023}, our models are free from relying on powerful close-sourced models like GPT-4 to improve the performance in low-resource languages. Our self-preferencing method also shares certain flavors with another self-rewarding mechanism \citep{yuan2024self_reward_lm}.\footnote{Our work was publicly available before \citet{yuan2024self_reward_lm}.}

\section{Evaluation}\label{sec:eval}

\subsection{Model Variants}\label{sec:eval:variants}

We trained multiple variants of SeaLLMs, as specified in the following.

\begin{itemize}
    \item \textbf{SeaLLM-7B-v1}: Trained from Llama-2-7B, it supports the 10 official languages used in Southeast Asia.
    \item \textbf{SeaLLM-13B-v1}: Trained from Llama-2-13B, it outperforms ChatGPT-3.5 in most non-Latin SEA languages (Khm, Lao, Mya and Tha) by large margins.
    \item \textbf{SeaLLM-7B-v2}: Trained from Mistral-7B, it outperforms SeaLLM-13B-v1 by far in higher-resource SEA languages (Vie, Ind, Tha), and surpasses ChatGPT-3.5 in math reasoning in SEA languages.
    \item \textbf{SeaLLM-7B-v2.5}: Trained from Gemma-7B, it outperforms SeaLLM-7B-v2 and SeaLLM-13B-v1 remarkably and surpasses ChatGPT-3.5 in various aspects in SEA languages, especially non-Latin languages.
\end{itemize}

\subsection{Sea-bench Peer Comparison}\label{sec:eval:peer_com}

While there are popular benchmarks to evaluate LLMs as a helpful assistant, such as MT-bench \citep{mt_bench_zheng2023judging}, they are only English-based and not suitable for evaluating performances in low-resource languages. Due to such a lack of multilingual benchmarks for assistant-style models, we engaged native linguists to build a multilingual test set with instructions that cover SEA languages, called \textbf{Sea-bench}. The linguists sourced such data by translating open-source English test sets, collecting real user questions from local forums and websites, collecting real math and reasoning questions from reputable sources, as well as writing test instructions themselves. Our Sea-Bench consists of diverse categories of instructions to evaluate the models, as follows:

\begin{itemize}
    \item Task-solving: This type of data comprises various text understanding and processing tasks, such as summarization, translation, etc.
    \item Math-reasoning: This includes math problems and logical reasoning tasks.
    \item General-instruction data: This consists of general user-centric instructions, which evaluate the model's ability in general knowledge and writing.
    \item NaturalQA: This consists of queries posted by real users, often in popular local forums, with a variety of subjects and topics of local interest. The aim is to test the model’s capacity to understand and respond coherently to colloquial language, natural expressions and idiomatic language, and locally contextualized references.
    \item Safety: This includes both general safety and local context-related safety instructions. While most general safety questions are translated from open sources, other local country-specific safety instructions are written by linguists of each language.
\end{itemize}

As inspired by MT-bench \citep{mt_bench_zheng2023judging}, we evaluate and compare SeaLLMs with well-known and state-of-the-art models using GPT-4 as a judge in a score-based grading metrics and a peer comparison (or pairwise comparison) manner. 

\ref{fig:sea_bench:seallm_vs_chatgpt} compares our SeaLLM (v2, v2.5) chat models with Qwen1.5-7B-chat \citep{qwen_bai2023} and the widely reputed ChatGPT-3.5\footnote{gpt-3.5-turbo June 2023 version.} \citep{chatgpt}. In the ``By Category'' chart, SeaLLM-7B-v2.5 performs on par with or surpasses ChatGPT-3.5 across various linguistic and writing tasks. This is largely thanks to the large gap in low-resource non-Latin languages, such as Burmese (Mya), Lao, Khmer and Thai, as seen in the ``By language'' chart on the right in \ref{fig:sea_bench:seallm_vs_chatgpt}. 

\subsection{MT-bench}\label{sec:eval:mt_bench}

We also compare our models with certain baselines on the English MT-Bench \citep{mt_bench_zheng2023judging} in \ref{table:eval:mt_bench}. As shown, SeaLLM-7B-v2 model demonstrates outstanding ability in English, given its size. It is also a rare multilingual model in the 7B realm, especially since it focuses on non-mainstream languages.

\subsection{World Knowledge}\label{sec:eval:knowledge}

In this section, we evaluate our models and reputable chat baselines \citep{llama2touvron2023llama,polylm_wei2023polylm,chatgpt} in terms of world knowledge. For knowledge across languages, we use the M3Exam benchmark \citep{m3exam}, which consists of real questions from human exam papers with various degrees of difficulty, ranging from primary school to high school examinations. We evaluate M3Exam with 3-shot native-instruction prompts across English, Chinese, Vietnamese, Indonesian and Thai. We also evaluate our models with the well-known English-centric MMLU benchmark \citep{mmlu_hendryckstest2021}.

\ref{table:eval:world_knowledge} details the evaluations of world knowledge across multiple languages and models of different sizes. SeaLLM-7B-v2.5 exhibits the best performance given its size and is competitive to GPT-3.5.

\subsection{Math Reasoning}\label{sec:eval:math_reasoning}

\ref{table:eval:sea_math} shows the GSM8K and MATH scores \citep{gsm8k_cobbe2021gsm8k,MATH_hendrycksmath2021} for zero-shot chain-of-thought prompting for English and their translated version in Chinese, Vietnamese, Indonesian and Thai. As shown, SeaLLM-7B-v2.5 shows competitive English performance in math reasoning compared to open-source models, with 78.5 in GSM8K and 34.9 in MATH. It also exceeds GPT-3.5 in SEA languages. This is achieved by scaling supervised and preference data in math reasoning in multilingual settings.

\subsection{Machine Translation}\label{sec:eval:translation}
To benchmark the machine translation performance of our SeaLLMs, we evaluate 4-shot chrF++ scores on the test sets from Flores-200~\cite{nllb_costa2022no_flores200}. As can be seen from \ref{fig:translation}, SeaLLM-13B exhibits clear superiority over ChatGPT-3.5 in low-resource languages, such as Lao and Khmer, while maintaining comparable performance with ChatGPT-3.5 in most higher resource languages (e.g., Vietnamese and Indonesian). 

For direct translation between SEA languages, as shown in \ref{fig:direct_translation}, our SeaLLM-13B-v1 model still achieves higher chrF++ scores than ChatGPT-3.5 in most cases, especially when the translation pairs involve low-resource languages. Overall, we believe our SeaLLMs will play a key role in facilitating communication and cultural exchange across communities in Southeast Asia.

\section{Conclusion}
In conclusion, our research presents a substantial advance in the development of equitable and culturally aware AI with the creation of SeaLLMs, a specialized suite of language models attuned to the linguistic and cultural landscapes of Southeast Asia. Through rigorous pre-training enhancements and culturally tailored fine-tuning processes, SeaLLMs have demonstrated exceptional proficiency in language understanding and generation tasks, challenging the performance of dominant players such as ChatGPT-3.5, particularly in SEA languages. The models' attunement to local norms and legal stipulations—validated by human evaluations—establishes SeaLLMs as not only a technical breakthrough but a socially responsive innovation, poised to democratize access to high-quality AI language tools across linguistically diverse regions. This work lays a foundation for further research into language models that respect and uphold the rich tapestry of human languages and cultures, ultimately driving the AI community towards a more inclusive future.

\section{Limitations}

SeaLLMs are among the most linguistically diverse multilingual large language models with remarkable abilities in languages beyond mainstream. However, they do not come without limitations. First, they only scratch the surface of the regionally linguistic diversity with 9 most common and representative languages, while there are hundreds other languages spoken in the Southeast Asia, such as Javanese and Tamil. Second, despite outperforming other popular models in non-Latin low-resource languages, SeaLLM models still suffer from considerable hallucination and degeneration under certain circumstances for languages such as Burmese and Lao. Moderate hallucination is still inevitable for other common languages.

\end{document}
